\documentclass[a4paper, serif, 10pt]{moderncv}

%% moderncv
% style and color
\moderncvstyle{banking}
\moderncvcolor{black}

%% page layout/margins
\usepackage[top=2.5cm, bottom=2.5cm, left=1.5cm, right=1.5cm]{geometry}

\nopagenumbers{}

%% Babel config for Indonesian (Bahasa Indonesia) language
% ref:
% - https://latex3.github.io/babel/guides/locale-indonesian.html
\usepackage[indonesian]{babel}

%% fontspec
\usepackage{fontspec}

\IfFontExistsTF{Linux Libertine}{
  \setmainfont[Ligatures={TeX, Common}, Numbers=OldStyle]{Linux Libertine}
}{}
\IfFontExistsTF{Linux Libertine O}{
  \setmainfont[Ligatures={TeX, Common}, Numbers=OldStyle]{Linux Libertine O}
}{}

%% CTeX
% CJK support, used to show CJK characters in the resume
%
% - fontset=none: disable builtin fontset but instead set the CJK font manually
% - heading=false: disable ctex heading
% - punct=kaiming: use kaiming punctuations styles for CJK
% - scheme=plain: use plain scheme, do not override \normalsize font size
% - space=auto: space settings for CJK characters
%
% ref:
% - http://ctan.mirrorcatalogs.com/language/chinese/ctex/ctex.pdf
\usepackage[UTF8, heading=false, punct=kaiming, scheme=plain, space=auto]{ctex}

\IfFontExistsTF{Noto Serif CJK SC}{
  \setCJKmainfont{Noto Serif CJK SC}
}{}
\IfFontExistsTF{Noto Sans CJK SC}{
  \setCJKsansfont{Noto Sans CJK SC}
}{}

%% line spacing
\usepackage{setspace}
\setstretch{1.125}

%% URL styling - use normal text instead of monospace
\urlstyle{same}

% Auto-underline all links
% Defer until \begin{document} so hyperref is already loaded
\AtBeginDocument{%
  \let\oldhref\href
  \renewcommand{\href}[2]{\oldhref{#1}{\underline{#2}}}%
}

%% Patch \httplink and \httpslink to handle full URLs
% The original moderncv macros always prepend http:// or https:// to URLs.
% This causes issues when the URL already has a protocol (e.g., https://example.com)
% resulting in malformed URLs like https://https://example.com.
% This patch checks if the URL already contains "://" and uses it directly if so.
% Note: We use \str_set:Nx (x-expansion) to expand macros like \@homepage first.
\ExplSyntaxOn
\str_new:N \g_yamlresume_url_str

\RenewDocumentCommand{\httpslink}{O{}m}{
  \str_set:Nx \g_yamlresume_url_str {#2}
  \str_if_in:NnTF \g_yamlresume_url_str {://}
    { \url{#2} }
    { \url{https://#2} }
}

\RenewDocumentCommand{\httplink}{O{}m}{
  \str_set:Nx \g_yamlresume_url_str {#2}
  \str_if_in:NnTF \g_yamlresume_url_str {://}
    { \url{#2} }
    { \url{http://#2} }
}
\ExplSyntaxOff

\name{Dewi Anggraini}{}
\title{Manajer Pemasaran dengan pengalaman 8+ tahun dalam strategi digital dan manajemen merek.}
\phone[mobile]{+62 812-3456-7890}
\email{dewi.anggraini@email.com}
\homepage{https://www.dewianggraini.com}

\address{Jl. Jenderal Sudirman No. 45, SCBD, Jakarta, DKI Jakarta, Indonesia, 12190}{}{}

\extrainfo{{\small \faLinkedin}\ \href{https://www.linkedin.com/in/dewianggraini}{@dewianggraini} {} {} {} • {} {} {} 
{\small \faTwitter}\ \href{https://twitter.com/dewianggraini}{@dewianggraini} {} {} {} • {} {} {} 
{\small \faInstagram}\ \href{https://www.instagram.com/dewianggraini}{@dewianggraini}}

\begin{document}

\maketitle

\section{Informasi Dasar}

\cvline{}{\begin{itemize}
\item Manajer Pemasaran dengan pengalaman lebih dari 8 tahun di industri ritel dan teknologi, terbukti meningkatkan pendapatan melalui strategi omni-channel.
\item Ahli dalam pemasaran digital, riset pasar, dan manajemen merek; memimpin tim lintas fungsi untuk kampanye yang berdampak.
\item Terbiasa bekerja dengan data untuk mengoptimalkan ROI, serta membangun hubungan kuat dengan mitra dan pelanggan.
\end{itemize}}
\section{Pendidikan}

\cventry{Sep 2010–Jul 2014}
        {Sarjana, Ilmu Komunikasi, Nilai: 3.75}
        {Universitas Indonesia}
        {\url{https://www.ui.ac.id}}
        {}
        {\begin{itemize}
\item Lulus dengan predikat cum laude, aktif dalam organisasi kemahasiswaan dan kompetisi pemasaran.
\item Mengembangkan keterampilan analitis dan kreatif melalui proyek kampanye sosial dan bisnis.
\item Menjadi asisten dosen untuk mata kuliah Riset Pemasaran selama satu semester.
\end{itemize}
\textbf{Mata Kuliah}: Manajemen Pemasaran, Perilaku Konsumen, Riset Pemasaran, Komunikasi Massa, Public Relations, Strategi Branding, Periklanan, Statistik Sosial}
\section{Pengalaman Kerja}

\cventry{Jan 2020–Sekarang}
        {Manajer Pemasaran}
        {PT Maju Bersama}
        {\url{https://www.majubersama.co.id}}
        {}
        {\begin{itemize}
\item Memimpin tim pemasaran 12 orang, mengelola anggaran tahunan Rp15 miliar, dan meningkatkan pendapatan 35\% YoY.
\item Merancang dan mengeksekusi strategi omni-channel (digital, offline, CRM) yang meningkatkan akuisisi pelanggan 40\%.
\item Mengembangkan kampanye brand awareness yang meraih 3 penghargaan industri dan meningkatkan brand recall 25\%.
\item Menganalisis data pasar dan kompetitor untuk mengidentifikasi peluang baru serta mengoptimalkan bauran pemasaran.
\item Membangun kemitraan strategis dengan influencer dan media, menghasilkan liputan senilai Rp2 miliar.
\end{itemize}
\textbf{Kata Kunci}: Strategi Pemasaran, Omni-channel, Manajemen Anggaran, Brand Awareness, CRM, Analitik Pemasaran, Kepemimpinan Tim}

\cventry{Mar 2016–Des 2019}
        {Asisten Manajer Pemasaran}
        {PT Kreatif Digital}
        {\url{https://www.kreatifdigital.com}}
        {}
        {\begin{itemize}
\item Mengelola kampanye digital untuk 20+ klien, dengan rata-rata peningkatan ROI 150\%.
\item Mengoptimalkan SEO/SEM dan media sosial, meningkatkan traffic organik 200\% dalam 18 bulan.
\item Menyusun laporan kinerja bulanan dan memberikan rekomendasi berbasis data kepada manajemen.
\item Berkolaborasi dengan tim kreatif dan teknologi untuk meluncurkan situs web dan aplikasi mobile.
\end{itemize}
\textbf{Kata Kunci}: Digital Marketing, SEO, SEM, Media Sosial, Google Analytics, Content Marketing, ROI}
\section{Bahasa}

\cvline{Indonesia}{Bahasa Ibu / Bilingual}
\cvline{Inggris}{Tingkat Profesional Penuh \hfill \textbf{Kata Kunci}: TOEFL 110, IELTS 8.0}
\section{Keahlian}

\cvline{Pemasaran Digital}{Ahli \hfill \textbf{Kata Kunci}: SEO/SEM, Media Sosial, Email Marketing, Content Strategy, Google Ads, Meta Ads, Analytics}
\cvline{Manajemen Merek}{Mahir \hfill \textbf{Kata Kunci}: Brand Strategy, Positioning, Consumer Insight, Kampanye Integratif, Public Relations}
\cvline{Analitik \& Riset}{Mahir \hfill \textbf{Kata Kunci}: Google Analytics, Tableau, SPSS, A/B Testing, Riset Pasar, Segmentasi}
\cvline{Kepemimpinan}{Ahli \hfill \textbf{Kata Kunci}: Manajemen Tim, Mentoring, Presentasi, Negosiasi, Manajemen Proyek}
\section{Penghargaan}

\cventry{Des 2022}
        {Asosiasi Pemasaran Indonesia}
        {Marketer of the Year 2022}
        {}
        {}
        {Diberikan kepada manajer pemasaran yang menunjukkan inovasi luar biasa dalam kampanye digital dan kontribusi signifikan terhadap pertumbuhan industri.}

\cventry{Nov 2021}
        {Digital Marketing Awards}
        {Best Digital Campaign}
        {}
        {}
        {Menghargai kampanye \#BanggaLokal yang berhasil meningkatkan keterlibatan pelanggan sebesar 300\% dan penjualan 50\% dalam 6 bulan.}
\section{Sertifikat}

\cventry{Mei 2022}
        {Google}
        {Google Digital Marketing \& E-commerce Professional Certificate}
        {\url{https://www.coursera.org/professional-certificates/google-digital-marketing-ecommerce}}
        {}
        {}

\cventry{Agu 2021}
        {HubSpot Academy}
        {HubSpot Content Marketing Certification}
        {\url{https://academy.hubspot.com/courses/content-marketing}}
        {}
        {}

\cventry{Feb 2023}
        {Meta}
        {Meta Certified Digital Marketing Associate}
        {\url{https://www.facebook.com/business/learn/certification}}
        {}
        {}
\section{Publikasi}

\cventry{Jun 2023}
        {Tren Pemasaran Digital di Indonesia 2023: Strategi Omni-channel untuk UMKM}
        {Jurnal Pemasaran Indonesia}
        {\url{https://www.jurnalpemasaran.id/artikel/tren-2023}}
        {}
        {\begin{itemize}
\item Menganalisis adopsi strategi omni-channel oleh UMKM di Indonesia.
\item Menyajikan studi kasus dari 50 bisnis yang berhasil meningkatkan pendapatan melalui integrasi digital dan offline.
\item Merekomendasikan kerangka kerja untuk pemasaran yang berkelanjutan dan berbasis data.
\end{itemize}}
\section{Proyek}

\cventry{Jan 2022–Des 2022}
        {Kampanye pemasaran terintegrasi untuk mendorong konsumsi produk lokal.}
        {Kampanye \#BanggaLokal}
        {\url{https://www.banggalokal.id}}
        {}
        {\begin{itemize}
\item Merancang dan memimpin kampanye multi-channel yang melibatkan 100+ influencer dan 50 media nasional.
\item Meningkatkan penjualan produk lokal sebesar 50\% dan engagement media sosial 300\%.
\item Meraih penghargaan Best Digital Campaign 2021 dan menjadi studi kasus di universitas.
\end{itemize}
\textbf{Kata Kunci}: Kampanye Integratif, Influencer Marketing, Media Sosial, Brand Lokal}

\cventry{Mar 2021–Sep 2021}
        {Peluncuran aplikasi belanja online untuk pasar Indonesia.}
        {Peluncuran Aplikasi Mobile BelanjaMudah}
        {\url{https://www.belanjamudah.id}}
        {}
        {\begin{itemize}
\item Mengoordinasikan tim pemasaran, produk, dan teknologi untuk peluncuran aplikasi.
\item Mencapai 100.000 unduhan dalam 3 bulan pertama melalui strategi user acquisition.
\item Mengelola anggaran peluncuran Rp5 miliar dengan ROI 200\%.
\end{itemize}
\textbf{Kata Kunci}: Product Launch, User Acquisition, Mobile Marketing, ROI}
\section{Minat}

\cvline{Fotografi}{Street Photography, Portrait, Landscape}
\cvline{Membaca}{Buku Bisnis, Fiksi, Sejarah}
\cvline{Olahraga}{Yoga, Berlari, Bersepeda}
\section{Kegiatan Sukarela}

\cventry{Feb 2019–Des 2021}
        {Relawan Pemasaran}
        {Yayasan Pendidikan Anak Indonesia}
        {\url{https://www.ypai.org}}
        {}
        {\begin{itemize}
\item Membantu mengembangkan strategi pemasaran digital untuk meningkatkan donasi dan kesadaran publik.
\item Mengelola kampanye media sosial yang berhasil mengumpulkan dana Rp500 juta untuk beasiswa anak kurang mampu.
\item Melatih staf yayasan dalam penggunaan alat digital untuk komunikasi dan penggalangan dana.
\end{itemize}}
    

\cventry{Jan 2018–Jan 2020}
        {Koordinator Acara}
        {Komunitas Pemasaran Jakarta}
        {\url{https://www.marketingjakarta.org}}
        {}
        {\begin{itemize}
\item Mengorganisir seminar dan workshop bulanan untuk 200+ profesional pemasaran.
\item Menghubungkan pembicara ahli dan sponsor untuk mendukung pengembangan karier anggota.
\end{itemize}}
    
\end{document}